\section{Conclusion}\label{sec:conclusion}

Firms occupy distribution-valued information positions rather than single
economic locations.
This paper turns target-anchored Wasserstein barycentric reconstruction into a
cross-sectional interaction field for those positions.
For each target, the fitted distributional spanning weights combine aligned peer
positions; repeated across firms, they form the admissible, directed barycentric
interaction field $W^\flat$, with nonnegative unit-sum rows and a zero diagonal.

Placing that field inside the quadratic exposure-adjustment problem yields
spatial closure: peer-adjusted exposure balances its stand-alone counterpart
against the peer average and obeys the familiar spatial autoregression.
Within this closure, $\rho=\lambda/(1+\lambda)$ records the relative intensity
of peer alignment, so spatial feedback acquires an exposure-adjustment
interpretation rather than remaining an unexplained coefficient conditional on
a supplied matrix.

In the 2023--2026 return panel, the barycentric interaction field (constructed
without using those evaluation returns) organizes conditional cross-sectional
return dependence and delivers stronger
conditional fit than pairwise RBF proximity constructed from the same
Wasserstein distances.
It selects a distinct peer structure and remains incrementally informative when
a persistent field of explicit news co-mention links enters the same
specification.
Joint estimation reallocates model-implied adjustment between these fields
rather than materially increasing total feedback.

Methodologically, the paper provides a bridge from distribution-valued
representation to spatial econometrics: it converts target-specific joint
representability in Wasserstein space into an admissible field for spatial
quasi-likelihood.
As Section~\ref{sec:identification-estimation} explains, the estimates are
conditional working-model quantities, not causal peer effects or
geometry-invariant structural adjustment parameters.
Spatial econometrics therefore need not begin after the researcher has supplied
$W$: it can begin one level upstream, by constructing the field from firms'
distribution-valued information positions.
